% GENERATED FILE. Edit canonical lessons, then run npm run generate:books.
% canonical-source-hash: fnv1a64:55caa6a04f6decb2
% canonical-lessons: JA-C382-kaigan, JA-C382-sunahama, JA-C382-misaki, JA-C382-suiheisen, JA-C382-kagayaku

\chapter{Down to the Coast}
\label{ch:ja-down-to-the-coast}
\hlchaptermodality{382}

\begin{chapteropening}
\textbf{By the end of this chapter:} \emph{I can describe a walk along the coast.}

Complete the last lesson of chapter 382: I can describe a walk along the coast.
\end{chapteropening}

\section[\texorpdfstring{\ja{かいがん} (kaigan)}{かいがん (kaigan)}]{\ja{かいがん} (kaigan) — the coast; the seashore}

\label{lesson:JA-C382-kaigan}



\begin{quote}
Before the new one: say the Japanese for to row, then the Japanese for to drift.
\end{quote}

\subsection*{You'll want to know: \ja{かいがん}}
\textbf{\ja{かいがん}} — \emph{kaigan} — \textquotedblleft{}the coast; the seashore\textquotedblright{}.

\begin{quote}
\textbf{\ja{なつに} \ja{かいがんを} \ja{あるきました。}} — \emph{natsu ni kaigan o arukimashita} — \textquotedblleft{}In summer I walked along the coast.\textquotedblright{}
\end{quote}

\begin{cousinweb}
\textbf{\ja{かいがん}} (\emph{kaigan}) is Sino-Japanese (kango): two Chinese characters read with their on'yomi, readings borrowed from Chinese. They are \emph{kai}, \textquotedblleft{}sea\textquotedblright{} (Middle Chinese \emph{xojX}) and \emph{gan}, \textquotedblleft{}beach\textquotedblright{} (\emph{nganH}) (KANJIDIC2; Middle Chinese in Baxter's transcription).
\end{cousinweb}

\begin{grammarlens}[title={What you've built}]
One more word for getting around.
\end{grammarlens}

\subsection*{Your turn}
\begin{itemize}
\raggedright
  \item \emph{Say it:} \emph{kaigan}
  \item \emph{Say it:} \emph{kaigan}, once more
  \item \emph{Say it:} \emph{natsu ni kaigan o arukimashita} — In summer I walked along the coast.
  \item \emph{Recall:} say the Japanese for to row, then the Japanese for to drift, then say \emph{kaigan} again
\end{itemize}

\begin{tcolorbox}[breakable,colback=teal!4,colframe=teal!35!black,title={Before you move on}]
What does \ja{かいがん} mean? (\textquotedblleft{}The coast; the seashore\textquotedblright{}.) Say it once more.
\end{tcolorbox}
\section[\texorpdfstring{\ja{すなはま} (sunahama)}{すなはま (sunahama)}]{\ja{すなはま} (sunahama) — a sandy beach}

\label{lesson:JA-C382-sunahama}



\begin{quote}
Before the new one: say the Japanese for to drift, then the Japanese for the coast.
\end{quote}

\subsection*{You'll want to know: \ja{すなはま}}
\textbf{\ja{すなはま}} — \emph{sunahama} — \textquotedblleft{}a sandy beach\textquotedblright{}.

\begin{quote}
\textbf{\ja{こどもが} \ja{すなはまで} \ja{あそんでいます。}} — \emph{kodomo ga sunahama de asonde imasu} — \textquotedblleft{}The children are playing on the sandy beach.\textquotedblright{}
\end{quote}

\begin{cousinweb}
\textbf{\ja{すなはま}} (\emph{sunahama}) is a native Japanese word (wago), inherited rather than borrowed from Chinese, built from \emph{suna}, \textquotedblleft{}sand\textquotedblright{} and \emph{hama}, \textquotedblleft{}beach\textquotedblright{}; its kanji were attached later as kun'yomi (JMdict).
\end{cousinweb}

\begin{grammarlens}[title={What you've built}]
One more word for getting around.
\end{grammarlens}

\subsection*{Your turn}
\begin{itemize}
\raggedright
  \item \emph{Say it:} \emph{sunahama}
  \item \emph{Say it:} \emph{sunahama}, once more
  \item \emph{Say it:} \emph{kodomo ga sunahama de asonde imasu} — The children are playing on the sandy beach.
  \item \emph{Recall:} say the Japanese for to drift, then the Japanese for the coast, then say \emph{sunahama} again
\end{itemize}

\begin{tcolorbox}[breakable,colback=teal!4,colframe=teal!35!black,title={Before you move on}]
What does \ja{すなはま} mean? (\textquotedblleft{}A sandy beach\textquotedblright{}.) Say it once more.
\end{tcolorbox}
\section[\texorpdfstring{\ja{みさき} (misaki)}{みさき (misaki)}]{\ja{みさき} (misaki) — a cape; a headland}

\label{lesson:JA-C382-misaki}



\begin{quote}
Before the new one: say the Japanese for the coast, then the Japanese for a sandy beach.
\end{quote}

\subsection*{You'll want to know: \ja{みさき}}
\textbf{\ja{みさき}} — \emph{misaki} — \textquotedblleft{}a cape; a headland\textquotedblright{}.

\begin{quote}
\textbf{\ja{みさきから} \ja{うみが} \ja{よく} \ja{みえます。}} — \emph{misaki kara umi ga yoku miemasu} — \textquotedblleft{}You can see the sea well from the cape.\textquotedblright{}
\end{quote}

\begin{cousinweb}
\textbf{\ja{みさき}} (\emph{misaki}) is a native Japanese word (wago), inherited rather than borrowed from Chinese; its kanji was attached later as a kun'yomi (JMdict).
\end{cousinweb}

\begin{grammarlens}[title={What you've built}]
One more word for getting around.
\end{grammarlens}

\subsection*{Your turn}
\begin{itemize}
\raggedright
  \item \emph{Say it:} \emph{misaki}
  \item \emph{Say it:} \emph{misaki}, once more
  \item \emph{Say it:} \emph{misaki kara umi ga yoku miemasu} — You can see the sea well from the cape.
  \item \emph{Recall:} say the Japanese for the coast, then the Japanese for a sandy beach, then say \emph{misaki} again
\end{itemize}

\begin{tcolorbox}[breakable,colback=teal!4,colframe=teal!35!black,title={Before you move on}]
What does \ja{みさき} mean? (\textquotedblleft{}A cape; a headland\textquotedblright{}.) Say it once more.
\end{tcolorbox}
\section[\texorpdfstring{\ja{すいへいせん} (suiheisen)}{すいへいせん (suiheisen)}]{\ja{すいへいせん} (suiheisen) — the horizon (over the sea)}

\label{lesson:JA-C382-suiheisen}



\begin{quote}
Before the new one: say the Japanese for a sandy beach, then the Japanese for a cape.
\end{quote}

\subsection*{You'll want to know: \ja{すいへいせん}}
\textbf{\ja{すいへいせん}} — \emph{suiheisen} — \textquotedblleft{}the horizon (over the sea)\textquotedblright{}.

\begin{quote}
\textbf{\ja{すいへいせんに} \ja{ゆうひが} \ja{しずみました。}} — \emph{suiheisen ni yūhi ga shizumimashita} — \textquotedblleft{}The evening sun sank below the horizon.\textquotedblright{}
\end{quote}

\begin{cousinweb}
\textbf{\ja{すいへいせん}} (\emph{suiheisen}) is Sino-Japanese (kango): three Chinese characters read with their on'yomi, readings borrowed from Chinese. They are \emph{sui}, \textquotedblleft{}water\textquotedblright{} (Middle Chinese \emph{sywijX}), \emph{hei}, \textquotedblleft{}flat\textquotedblright{} (\emph{bjaeng}) and \emph{sen}, \textquotedblleft{}line\textquotedblright{} (\emph{sjenH}) (KANJIDIC2; Middle Chinese in Baxter's transcription).
\end{cousinweb}

\begin{grammarlens}[title={What you've built}]
One more word for getting around.
\end{grammarlens}

\subsection*{Your turn}
\begin{itemize}
\raggedright
  \item \emph{Say it:} \emph{suiheisen}
  \item \emph{Say it:} \emph{suiheisen}, once more
  \item \emph{Say it:} \emph{suiheisen ni yūhi ga shizumimashita} — The evening sun sank below the horizon.
  \item \emph{Recall:} say the Japanese for a sandy beach, then the Japanese for a cape, then say \emph{suiheisen} again
\end{itemize}

\begin{tcolorbox}[breakable,colback=teal!4,colframe=teal!35!black,title={Before you move on}]
What does \ja{すいへいせん} mean? (\textquotedblleft{}The horizon (over the sea)\textquotedblright{}.) Say it once more.
\end{tcolorbox}
\section[\texorpdfstring{\ja{かがやく} (kagayaku)}{かがやく (kagayaku)}]{\ja{かがやく} (kagayaku) — to sparkle; to shine brightly}

\label{lesson:JA-C382-kagayaku}



\begin{quote}
Before the new one: say the Japanese for a cape, then the Japanese for the horizon (over the sea).
\end{quote}

\subsection*{You'll want to know: \ja{かがやく}}
\textbf{\ja{かがやく}} — \emph{kagayaku} — \textquotedblleft{}to sparkle; to shine brightly\textquotedblright{}.

\begin{quote}
\textbf{\ja{うみが} \ja{あさひで} \ja{かがやいています。}} — \emph{umi ga asahi de kagayaite imasu} — \textquotedblleft{}The sea is sparkling in the morning sun.\textquotedblright{}
\end{quote}

Polite \textbf{\ja{かがやきます}}, \emph{kagayakimasu}; a godan verb. Past (polite): \textbf{\ja{うみが} \ja{かがやきました。}}, \emph{umi ga kagayakimashita}.

\begin{cousinweb}
\textbf{\ja{かがやく}} (\emph{kagayaku}) is a native Japanese word (wago), inherited rather than borrowed from Chinese; its kanji was attached later as a kun'yomi (JMdict).
\end{cousinweb}

\begin{grammarlens}[title={What you've built}]
One more word for getting around.
\end{grammarlens}

\subsection*{Your turn}
\begin{itemize}
\raggedright
  \item \emph{Say it:} \emph{kagayaku}
  \item \emph{Say it:} \emph{kagayaku}, once more
  \item \emph{Say it:} \emph{umi ga asahi de kagayaite imasu} — The sea is sparkling in the morning sun.
  \item \emph{Recall:} say the Japanese for a cape, then the Japanese for the horizon (over the sea), then say \emph{kagayaku} again
\end{itemize}

\begin{tcolorbox}[breakable,colback=teal!4,colframe=teal!35!black,title={Before you move on}]
What does \ja{かがやく} mean? (\textquotedblleft{}To sparkle; to shine brightly\textquotedblright{}.) Say it once more.
\end{tcolorbox}
